% Solutions/hints for ch11 exercises (assembled into Appendix C)

\begin{solution}{exr:ch11-sizes}
(a) With $|V| = 300$ there are $300^3 = 2.7 \cdot 10^7$ joint configurations for three agents, of which $300 \cdot 299 \cdot 298 \approx 2.67 \cdot 10^7$ are collision-free, and $5^3 = 125$ joint moves per node. For twelve agents there are $300^{12} \approx 5.3 \cdot 10^{29}$ configurations, about $300^{\underline{12}} \approx 4.2 \cdot 10^{29}$ collision-free ones (the product $300 \cdot 299 \cdots 289$), and $5^{12} \approx 2.4 \cdot 10^8$ joint moves per node. (b) One full joint step of plain joint \astar generates all $5^{12} \approx 2.4 \cdot 10^8$ successors of a node. Operator decomposition generates $5$ successors per intermediate state and needs $12$ levels; if nothing were pruned, the twelve levels would again enumerate $5^{12}$ leaves, but every level is an \astar node with its own $\fcost$, so the search expands only the intermediate states whose $\fcost$ is small enough, and in practice a handful of states per level: the enumeration is lazy and pruned by the heuristic, which the Cartesian product of plain joint \astar cannot be. (c) A heuristic decides which node is expanded next, not how many successors an expansion generates; every expansion of plain joint \astar enumerates $(b+1)^k$ joint moves and checks each for collisions, so even a search that expands only the nodes on an optimal plan pays $(b+1)^k$ per expansion.
\end{solution}

\begin{solution}{exr:ch11-collision-sets}
The policies are $\phi_1(a) = b$, $\phi_1(b) = c$, $\phi_2(c) = b$, $\phi_2(b) = a$ and $\phi_3(d) = d$ (agent~3 is at its goal and rests for free). Step~1 expands the start $(a, c, d)$ with $C = \emptyset$ and generates the policy successor $(b, b, d)$, in which agents~1 and~2 share the vertex $b$: the successor is discarded, $\set{1,2}$ is stored on it and backpropagated to the start, whose collision set becomes $\set{1,2}$; the start is re-opened. There is no node expanded with an empty collision set other than the start, so step~2 is the re-expansion of the start with $C = \set{1,2}$: agent~1 chooses among $\set{a, b}$, agent~2 among $\set{c, b}$ and agent~3 follows its policy and stays at $d$, which gives the limited neighbors $(a, c, d)$ (the joint wait, cost $2$), $(a, b, d)$, $(b, c, d)$ and $(b, b, d)$ (the known collision). Two of them enter \Open with $\gcost = 2$ and $\fcost = 5$: $(a, b, d)$ and $(b, c, d)$ (one agent moves, the other waits away from its goal, agent~3 rests at its goal for free). The joint wait leads back to the start itself, whose $\gcost$ is already $0$, so it is generated but never queued---a useful reminder that a limited neighbor is not automatically a new node. Step~3 expands one of the two nodes with $\fcost = 5$, whichever was pushed first; its collision set is empty, and its policy successor is a swap of agents~1 and~2 along the edge $a$--$b$ or $b$--$c$, so $\set{1,2}$ is backpropagated to it and to the start, and the node is re-opened. Agent~3 enters a collision set at expansion~5, the second expansion of $(a, b, d)$. Once agent~2 is coupled it may leave its policy, and one of its limited neighbors there is the move $b \to d$ into the vertex on which agent~3 rests: a vertex collision of agents~2 and~3. That set is stored on the successor and backpropagated, so $C((a, b, d))$ grows to $\set{1,2,3}$ and the start is re-opened with $C = \set{1,2,3}$ at expansion~6. The start therefore ends with the collision set $\set{1,2,3}$ and the search degenerates into full joint \astar over the three agents. The two-agent plan of cost $7$ needs $d$ as a passing place, and with agent~3 parked there the instance is unsolvable: agents~1 and~2 must pass each other on the path $a$--$b$--$c$ and cannot. \mstar therefore returns \emph{failure}, after $12$ expansions and $84$ generated candidates with the chapter code, and \cref{thm:ch11-terminates} guarantees that it returns at all, once every joint move of the three agents has been tried.
\end{solution}

\begin{solution}{exr:ch11-heuristic}
(a) Every agent $i$ must make at least $d_i(v^i)$ moves, each of cost one, so the cost to go is at least $\max_i d_i(v^i)$, and $\hcost_{\max} \le \hcost^*$. Since all $d_i \ge 0$, $\max_i d_i(v^i) \le \sum_i d_i(v^i)$, so the sum dominates the maximum; by the dominance theorem of \cref{ch:ch04} it expands no more nodes with $\fcost < C^*$. (b) If waiting were free everywhere, a wait of an agent away from its goal would cost $0$ while $d_i$ stays the same, so $\hcost(v) \le c(v, v') + \hcost(v')$ still holds (both sides are equal) and the heuristic remains consistent; but the joint graph now has zero-cost cycles (everybody waits), so $\gcost$ can no longer distinguish waiting from not waiting, the cheapest plan may contain arbitrarily many waits, and the search stays finite only because \astar never re-expands a node with the same $\gcost$; the returned plan minimizes the number of moves, which is a different objective. (c) For the makespan, let every joint move cost $1$ unless every agent rests at its goal, in which case it costs $0$; the admissible heuristic is $\hcost_{\max}(v) = \max_i d_i(v^i)$, because at least that many time steps remain. Along the policy path every step costs $1$ and reduces $\hcost_{\max}$ by exactly $1$ until the slowest agent arrives, so $\fcost$ is constant and the policy path from $v$ costs $\hcost_{\max}(v)$, a lower bound; hence it is an optimal continuation whenever it is conflict-free.
\end{solution}

\begin{solution}{exr:ch11-reexpansion}
(a) Seven of the twelve are re-expansions after a collision among the node's \emph{own} successors: the node was expanded, one of its limited neighbors turned out to be a vertex or a swap collision, and line~\ref{alg:ch11-mstar:bpcol} added the two agents to the node itself before backpropagating. The other five received their set from a child (line~\ref{alg:ch11-mstar:bpchild}): the collision was found one or more steps deeper and traveled up the back sets. The start is one of the five, which is why its two expansions are steps~1 and~4: at step~1 its single policy successor is legal, and the head-on meeting appears only when that successor is expanded at step~2. (b) The collision set of $v$ is defined by what the search finds \emph{below} $v$, and \mstar has no oracle for it: computing it in advance would mean knowing which agents interact on the way to the goal, which is the problem itself. So every node starts with $C(v) = \emptyset$, is expanded along the policies, and is re-opened when a descendant proves that this was too narrow; the start is expanded once with $C = \emptyset$ (step~1) and again with $C = \set{1,2}$ (step~4), and no rule that keeps optimality can avoid the first of the two. The price is bounded: a node is re-expanded at most once per strict growth of its collision set, hence at most $k$ times, which is the argument of \cref{thm:ch11-terminates}. (c) The instance of \cref{exr:ch11-collision-sets} does exactly this: on the path $a$--$b$--$c$ with the spur $b$--$d$, agents~1 and~2 must exchange places and agent~3 rests on $d$. The start is re-opened with $C = \set{1,2}$ at expansion~2, and once agent~2 is free to leave its policy it tries the move $b \to d$ onto agent~3, so the start is re-opened again with $C = \set{1,2,3}$ at expansion~6; the run ends in failure after $12$ expansions. A three-drone aisle with a single bay behaves the same way as soon as the third drone is parked in the bay.
\end{solution}

\begin{solution}{exr:ch11-swap-count}
(a) The leader is $d$ steps from the junction $v$ and the follower one step behind it. The multipush moves the leader $d$ times and, after each of its moves, the follower once into the vertex just vacated, which is $2d$ moves; at $v$ the leader stands on the junction and the follower on the vertex $u$ before it. The exchange is the fixed sequence of \cref{fig:ch11-push-swap}(b): six moves. The reversal undoes the $2d$ preparatory moves one by one, whoever stands on the target of each move stepping back to its source, which is again $2d$ moves. In total $2d + 6 + 2d = 4d + 6$, and nothing else moved because no other agent was cleared. (b) In the swap example the code executes $14$ moves: one push of $b$ from $c_1$ to $c_2$, one move of $a$ to its goal $c_1$, the swap, in which $b$ leads from $c_2$ to the junction $c_3$ with $d = 1$ and which therefore costs $4 \cdot 1 + 6 = 10$ moves, one move of $b$ to $c_0$ and one move of the displaced $a$ back to $c_1$. With the junction $d$ steps from $c_2$ the count is $4d + 10$, the makespan of the one-move-per-step plan is $4d + 10$, and its sum of costs is the sum of the times of the last moves of the two agents, $(4d + 10) + (4d + 9) = 8d + 19$, which is $27$ for $d = 1$. The optimal joint plan sends both agents toward the junction together, parks $a$ in the spur for one step while $b$ turns around behind it, and brings both back: with the junction three steps from $a$'s start, as in the example, that is makespan $7$ and sum of costs $14$, and in general each agent needs the walk to the junction and back plus one or two steps, so both quantities grow linearly in the distance as well, but with the slope of two walks instead of four and with both agents moving at the same time. (c) Number the six moves $a: v \to n_1$, $b: u \to v$, $b: v \to n_2$, $a: n_1 \to v$, $a: v \to u$, $b: n_2 \to v$. They pair into the three time steps $\set{1,2}$, $\set{3,4}$ and $\set{5,6}$: each of the three pairs is a following move (one agent enters the vertex the other leaves in the same step, which \cref{sec:ch11-problem} allows), the two moves of a pair run on different edges, and neither is a swap along one edge, so no pair has a conflict. No fewer than three steps are possible because each of the two agents has to make three moves and an agent moves at most once per step; three steps are therefore exactly optimal. In the multipush the follower always enters the vertex the leader is leaving, which is a following move, so all $d$ pairs of moves can run in $d$ time steps.
\end{solution}
